\documentclass[11pt,a4paper]{article}

\usepackage[utf8]{inputenc}
\usepackage[T1]{fontenc}
\usepackage[french,english]{babel}

\usepackage{amsmath,amssymb,amsthm}
\usepackage{mathrsfs}
\usepackage{graphicx}
\usepackage{hyperref}
\usepackage{geometry}
\usepackage{physics}
\usepackage{enumitem}
\setlist[itemize,1]{label=---}
\geometry{margin=2.5cm}
\usepackage{csquotes}

\numberwithin{equation}{section}

% Theorem environments
\newtheorem{theorem}{Theorem}[section]
\newtheorem{lemma}[theorem]{Lemma}
\newtheorem{proposition}[theorem]{Proposition}
\newtheorem{corollary}[theorem]{Corollary}

\theoremstyle{definition}
\newtheorem{definition}[theorem]{Definition}
\newtheorem{axiom}[theorem]{Axiom}
\newtheorem{remark}[theorem]{Remark}

\title{%
Gleason-Type Uniqueness of Born's Rule\\
for Spin-$\tfrac{1}{2}$ in the PDL Framework
}

\author{C\'edric Laubscher}

\date{\today}

\begin{document}

\selectlanguage{english} % main writing language

\maketitle

\begin{abstract}
We construct an operational axiomatisation of spin-$\tfrac12$ measurements adapted to the Projective Dynamic Logo (PDL) framework and establish a Gleason-type uniqueness theorem for the corresponding probability assignments. Imposing physically motivated conditions of normalisation, global phase invariance, rotational covariance, branch symmetry, regularity, and measurement repeatability, we demonstrate that the probability $P_+(\theta;a,b)$ of obtaining the ``$+$'' outcome when measuring spin along an axis with polar angle $\theta$ in the state $|\psi\rangle = a\,|+\rangle_z + b\,|-\rangle_z$ is necessarily given by Born’s rule. Specifically, every admissible probability assignment takes the form $P_+(\theta;a,b)=|a\cos(\theta/2)+b\sin(\theta/2)|^2$, with $P_-=1-P_+$.

We subsequently reanalyse the mixed-triangle measurement scheme previously formulated within PDL, in which a minimal $(4,6)$ closure encodes a spin doublet and a Stern–Gerlach-type device is represented by a finite active surface composed of mixed triangles. We prove that, under mild symmetry constraints on the apparatus, the effective probabilities generated by this scheme satisfy our operational axioms and thus fall within the scope of the uniqueness theorem. It follows that Born’s rule for spin-$\tfrac12$ measurements is not merely implementable in PDL, but is uniquely enforced by the stipulated operational and combinatorial conditions. This result provides an initial link between Gleason-type derivations and the discrete relational substrate underlying PDL, and suggests the feasibility of a more general uniqueness programme for quantum probabilities within this framework.
\end{abstract}

\section{Introduction}
\label{sec:introduction}

Gleason's theorem and its generalisations occupy a central position in the foundations of quantum theory. They establish that, under appropriate assumptions of non-contextuality and additivity, any probability measure defined on the projection lattice of a Hilbert space of dimension at least three must be represented by the Born rule.\cite{Gleason1957,BuschEtAlBook} In this sense, these results furnish a strong form of \emph{uniqueness} for quantum probability assignments, tightly constraining them through the structure of Hilbert space.

The Projective Dynamic Logo (PDL) framework pursues a distinct route towards quantum phenomenology. Rather than postulating an underlying Hilbert space, it represents physical systems as finite signed graphs selected by a coherence-optimisation principle.\cite{LaubscherIntroPDL2026,LaubscherLDP2026} Within this framework, a minimal $(4,6)$ closure has been identified as a logical prototype for a two-level system, and a combinatorial scheme based on mixed triangles and active surfaces has recently been shown to reproduce the Born rule for spin-$\tfrac12$ measurements in a Stern--Gerlach-type scenario.\cite{LaubscherBornPDL2026} The conceptual status of this result has, however, remained essentially \emph{propositional}: the construction establishes that the Born rule is \emph{implementable} within PDL, but it does not address the question of its \emph{uniqueness}.

The objective of the present work is to address this gap, at least at the level of spin-$\tfrac12$ systems. We formulate an operational axiomatisation of two-outcome spin measurements adapted to the PDL formalism and derive a first Gleason-type uniqueness theorem: under natural assumptions of normalisation, global phase invariance, rotational covariance, branch symmetry, regularity, and measurement repeatability, any admissible probability assignment $P_\pm(\theta;a,b)$ must coincide with the Born rule. We then demonstrate that the mixed-triangle measurement scheme of Ref.~\cite{LaubscherBornPDL2026} satisfies these axioms under mild symmetry requirements imposed on the PDL apparatus. Consequently, the Born rule is not merely one admissible choice within PDL, but rather the \emph{unique} probability assignment compatible with the stipulated operational and combinatorial constraints in the spin-$\tfrac12$ model.

Beyond its specific relevance for the PDL research programme, this result more broadly illustrates how Gleason-type uniqueness arguments can be reformulated in terms of a discrete relational substrate instead of the conventional Hilbert-space formalism. It thereby constitutes a first step towards a systematic treatment of quantum probabilities and measurement theory in PDL, and it indicates several promising directions for extending the approach to higher-dimensional systems and more general classes of observables.

\section[Operational setting for spin-1/2]{Operational setting for spin-$\tfrac{1}{2}$}
\label{sec:operational-setting}

In this section we formalize the operational framework for spin-$\tfrac12$
measurements that will serve as the basis for our subsequent uniqueness analysis.

\subsection{States and measurement directions}

We consider a two-dimensional quantum system prepared in a normalized pure state
\begin{equation}
|\psi\rangle = a\,|+\rangle_z + b\,|-\rangle_z,
\qquad a,b\in\mathbb{C},\quad |a|^2+|b|^2=1,
\end{equation}
where $\{|+\rangle_z,|-\rangle_z\}$ denotes a fixed reference eigenbasis of the spin observable along the $z$-axis. Measurement directions are parametrized by a polar angle $\theta\in\mathbb{R}$ with respect to the $z$-axis, restricted to a fixed plane.

We denote by $P_\pm(\theta;a,b)$ the operationally defined probabilities of obtaining the outcomes $+$ and $-$, respectively, when measuring the spin component along the axis forming an angle $\theta$ with the $z$-axis on the state $|\psi\rangle$.

\subsection{Operational axioms}
\label{subsec:axioms}

We now state the axioms that constrain the admissible probability assignments
\[
(\theta,a,b)\longmapsto P_\pm(\theta;a,b).
\]

\begin{axiom}[Normalization and positivity]
\label{ax:normalization}
For all $\theta\in\mathbb{R}$ and all normalized pairs $(a,b)$, the probabilities satisfy
\[
P_+(\theta;a,b)\ge 0,\quad P_-(\theta;a,b)\ge 0,\quad
P_+(\theta;a,b)+P_-(\theta;a,b)=1.
\]
\end{axiom}

\begin{axiom}[Global phase invariance]
\label{ax:global-phase}
For all $\theta\in\mathbb{R}$, all $(a,b)$ with $|a|^2+|b|^2=1$, and all
$\varphi\in\mathbb{R}$,
\[
P_{\pm}(\theta; e^{i\varphi}a, e^{i\varphi}b) = P_{\pm}(\theta;a,b).
\]
That is, the probabilities are invariant under multiplication of the state vector by an overall phase factor.
\end{axiom}

\begin{axiom}[Rotational covariance]
\label{ax:rotation}
For each real angle $\theta_0$ there exists a map
\[
R_{\theta_0} : (a,b)\mapsto(a',b')
\]
such that
\[
P_{\pm}(\theta;a',b') = P_{\pm}(\theta+\theta_0; a,b)
\]
for all $\theta\in\mathbb{R}$ and all normalized $(a,b)$. In other words, a rotation of the measurement direction by $\theta_0$ can be equivalently represented by a transformation $R_{\theta_0}$ acting on the state amplitudes.
\end{axiom}

\begin{axiom}[Eigenstate behavior at the reference angle]
\label{ax:eigenstates}
At the reference measurement direction $\theta=0$ (aligned with the $z$-axis) one has
\[
P_+(0;1,0)=1,\quad P_-(0;1,0)=0,\qquad
P_-(0;0,1)=1,\quad P_+(0;0,1)=0.
\]
Thus, the basis states $|+\rangle_z$ and $|-\rangle_z$ are eigenstates of the $\theta=0$ measurement with definite outcomes $+$ and $-$, respectively.
\end{axiom}

\begin{axiom}[Branch symmetry under $\theta\mapsto\theta+\pi$]
\label{ax:branch-symmetry}
For all $\theta\in\mathbb{R}$ and all normalized $(a,b)$, the probabilities obey
\[
P_+(\theta+\pi;a,b) = P_-(\theta;a,b),\qquad
P_-(\theta+\pi;a,b) = P_+(\theta;a,b).
\]
That is, a rotation of the measurement direction by $\pi$ reverses the identification of the two outcome labels.
\end{axiom}

\begin{axiom}[Regularity]
\label{ax:regularity}
The functions
\[
(\theta,a,b)\longmapsto P_{\pm}(\theta;a,b)
\]
are continuous in all their arguments and, for fixed $\theta$, analytic in $(a,b)$ when restricted to the unit sphere $|a|^2+|b|^2=1$ in $\mathbb{C}^2$.
\end{axiom}

\begin{axiom}[Repeatability in the eigenbasis]
\label{ax:repeatability}
Consider a measurement at the reference angle $\theta=0$. If the system is prepared in a state $|\psi\rangle$ such that $P_+(0;a,b)=1$ (respectively, $P_-(0;a,b)=1$), then an immediate repetition of the same measurement (performed on the post-measurement ensemble) yields the outcome $+$ (respectively, $-$) with probability $1$.

More generally, for any normalized $(a,b)$ and any integer $n\ge 1$, if one performs $n$ independent repetitions of the $\theta=0$ measurement on identically prepared copies of $|\psi\rangle$, the relative frequency of the outcome $+$ converges almost surely to $P_+(0;a,b)$ in the limit $n\to\infty$.
\end{axiom}

\section{A reduced-form lemma for spin-\texorpdfstring{$\tfrac12$}{1/2} probabilities}
\label{sec:reduced-form}

In this section we derive an initial structural constraint on the class of admissible probability assignments $P_\pm(\theta;a,b)$ satisfying Axioms
\ref{ax:normalization}--\ref{ax:regularity}.

% >>> Insert Lemma 1 here:

\begin{lemma}[Reduced form of spin-$\tfrac12$ probabilities]
\label{lem:reduced-form}
Consider a spin-$\tfrac12$ system prepared in a normalized state
\[
|\psi\rangle = a\,|+\rangle_z + b\,|-\rangle_z, \qquad a,b \in \mathbb{C},\quad |a|^2 + |b|^2 = 1,
\]
and let $P_\pm(\theta;a,b)$ denote the operational probabilities of obtaining the outcomes
$+$ or $-$ when measuring the spin component along an axis characterized by the polar angle $\theta$ with respect to the $z$-axis.

Assume that the probability functional $P_\pm$ satisfies axioms (A1)--(A6): normalization and positivity, invariance under global phase transformations, covariance with respect to rotations around the reference axis, correct limiting behavior for eigenstates at $\theta=0$, symmetry between outcome branches under the transformation $\theta \mapsto \theta + \pi$, and regularity (continuity/analyticity).

Then there exists a continuous function
\[
F : [0,1] \to [0,1]
\]
such that, for all normalized pairs $(a,b)$ and all $\theta \in \mathbb{R}$,
\begin{equation}
\label{eq:reduced-form}
P_+(\theta;a,b)
=
F\!\left(
\left|\,a\cos\frac{\theta}{2} + b\sin\frac{\theta}{2}\,\right|^2
\right),
\end{equation}
and
\begin{equation}
\label{eq:F-symmetry}
F(1 - x) = 1 - F(x), \qquad \forall\,x \in [0,1].
\end{equation}
In particular,
\[
P_-(\theta;a,b) = 1 - P_+(\theta;a,b)
= F\!\left(
\left|\,a\sin\frac{\theta}{2} - b\cos\frac{\theta}{2}\,\right|^2
\right).
\]
\end{lemma}

\begin{proof}
We now present a more detailed argument.

\medskip\noindent
\emph{Step 1: Reduced form at the reference angle.}
By Lemma~\ref{lem:reduced-form}, the probability of obtaining the outcome $+$ at the reference angle $\theta=0$ can be expressed as
\begin{equation}
\label{eq:P0-Fx}
P_+(0;a,b) = F(|a|^2),
\end{equation}
for some continuous function $F:[0,1]\to[0,1]$. In particular, for any normalized pair $(a,b)$, the single real parameter
\[
x := |a|^2 = |\langle +|\psi\rangle|^2
\]
completely determines $P_+(0;a,b)$.

\medskip\noindent
\emph{Step 2: Classical mixtures and convexity.}
Consider two normalized pure states $|\psi_1\rangle$, $|\psi_2\rangle$
with squared overlaps
\[
x_1 := |\langle +|\psi_1\rangle|^2,\qquad
x_2 := |\langle +|\psi_2\rangle|^2.
\]
For each $k\in\{1,2\}$, we have
\[
P_+^{(k)}(0) = F(x_k)
\]
for the measurement at $\theta=0$ performed on $|\psi_k\rangle$.

Now prepare a \emph{classical mixture} of these states: with probability $p\in[0,1]$ the preparation device outputs $|\psi_1\rangle$, and with probability $(1-p)$ it outputs $|\psi_2\rangle$, without disclosing to the measuring apparatus which state has been selected. By construction, the overall probability of obtaining the outcome $+$ in a single run of the
$\theta=0$ measurement on this ensemble is
\begin{equation}
\label{eq:mixture-prob}
P_+^{\mathrm{mix}}(0) = p\,F(x_1) + (1-p)\,F(x_2).
\end{equation}
This follows directly from the standard law of total probability applied to the
classical randomization occurring at the preparation stage.

\medskip\noindent
\emph{Step 3: A pure state with the same overlap.}
On the other hand, there always exists a normalized pure state
$|\psi\rangle$ whose squared overlap with $|+\rangle_z$ equals the convex
combination
\[
x := p x_1 + (1-p)x_2.
\]
Indeed, by choosing an appropriate vector in the two-dimensional subspace spanned by $|\psi_1\rangle$ and $|\psi_2\rangle$, one can realize any prescribed value $x\in[0,1]$ for $|\langle +|\psi\rangle|^2$.
For this state, Eq.~\eqref{eq:P0-Fx} implies
\begin{equation}
\label{eq:pure-prob}
P_+(0;\psi) = F(x) = F\bigl(p x_1 + (1-p)x_2\bigr).
\end{equation}

\medskip\noindent
\emph{Step 4: Identification via repeatability.}
Consider the following two experimental scenarios:

\begin{itemize}
\item[(E1)] One repeatedly measures the $\theta=0$ observable on independent copies of the \emph{classical mixture} of
$|\psi_1\rangle$ and $|\psi_2\rangle$ with classical weights $p$ and
$(1-p)$, respectively.
\item[(E2)] One repeatedly measures the same observable on independent copies of a \emph{pure state} $|\psi\rangle$ whose overlap with
$|+\rangle$ satisfies $|\langle +|\psi\rangle|^2 = x$.
\end{itemize}

By construction, in experiment (E1) the probability of obtaining the outcome $+$ in a single trial is $P_+^{\mathrm{mix}}(0)$, as defined in
\eqref{eq:mixture-prob}. In experiment (E2), the corresponding probability is $P_+(0;\psi)$, as given by \eqref{eq:pure-prob}. Applying the repeatability Axiom~\ref{ax:repeatability} to each experiment independently, and invoking the law of large numbers for independent and identically distributed trials, we obtain:
\begin{itemize}
\item In (E1), the empirical relative frequency of the outcome $+$ converges almost surely to $P_+^{\mathrm{mix}}(0)$.
\item In (E2), the empirical relative frequency of the outcome $+$ converges almost surely to $P_+(0;\psi)=F(x)$.
\end{itemize}

Assume now that the two experiments are performed using identical measurement devices, and that the only operationally relevant parameter at the reference angle $\theta=0$ is the single real number $x$, i.e.\ for any preparation, the probability of obtaining the outcome $+$ is uniquely determined by the value of $x$. In particular, an experimentalist who has access solely to the long-run statistics of the $\theta=0$ measurement cannot operationally discriminate whether a given value of $x$ is generated by a classical mixture or by a pure-state preparation. Consistency of the probability assignment therefore requires that experiments (E1) and (E2), which share the same effective value $x = p x_1 + (1-p)x_2$, yield identical limiting frequencies for the outcome $+$. Consequently,
\[
F\bigl(p x_1 + (1-p)x_2\bigr)
=
P_+(0;\psi)
=
P_+^{\mathrm{mix}}(0)
=
p\,F(x_1) + (1-p)\,F(x_2),
\]
for all $x_1,x_2\in[0,1]$ and all $p\in[0,1]$.

Equivalently, $F$ satisfies
\begin{equation}
\label{eq:F-Jensen}
F\bigl(p x_1 + (1-p)x_2\bigr)
=
p\,F(x_1) + (1-p)\,F(x_2),
\qquad
\forall\,x_1,x_2\in[0,1],\ \forall\,p\in[0,1].
\end{equation}
This is precisely Jensen's functional equation for convex combinations
on the interval $[0,1]$.

\medskip\noindent
\emph{Step 5: Affineness and symmetry.}
It is a standard result in the theory of functional equations that any continuous function $F:[0,1]\to\mathbb{R}$ satisfying
\eqref{eq:F-Jensen} must be affine. That is, there exist real constants
$\alpha,\beta\in\mathbb{R}$ such that
\[
F(x)=\alpha x + \beta,\qquad x\in[0,1].
\]
This establishes \eqref{eq:F-affine}.

Substituting this affine form into the symmetry condition
$F(1-x)=1-F(x)$ from \eqref{eq:F-symmetry} yields
\[
\alpha(1-x)+\beta = 1 - (\alpha x + \beta),
\]
for all $x\in[0,1]$. Identifying coefficients of $x$ and of the
constant term on both sides, we obtain
\[
\alpha = 1,\qquad \beta = 0,
\]
so that $F(x)=x$ on $[0,1]$, and thus \eqref{eq:F-identity} follows. The Born rule expressions for $P_\pm(\theta;a,b)$ then follow directly from Lemma~\ref{lem:reduced-form}.
\end{proof}

\begin{remark}
The pivotal aspect of the foregoing argument is that, at the reference setting $\theta = 0$, the probability of obtaining the outcome $+$ depends on the prepared quantum state solely through the scalar parameter $x = |\langle + | \psi \rangle|^2$. As a consequence, classical probabilistic mixtures and pure quantum superpositions that are characterized by the same value of $x$ cannot be operationally distinguished on the basis of their measurement outcome statistics. This constitutes precisely the form of operational non-contextuality that underlies Gleason-type derivations of Born’s rule.
\end{remark}

\section{Towards a Uniqueness Result for Born's Rule}
\label{sec:uniqueness}

\subsection{Initial Constraints on the Scalar Function \texorpdfstring{$F$}{F}}

In this subsection, we derive an initial structural constraint on the scalar function
$F:[0,1]\to[0,1]$ introduced in Lemma~\ref{lem:reduced-form}. We demonstrate that, under a mild additional assumption of measurement repeatability, the function $F$ must be affine.

\begin{proposition}[Affineness of $F$ under Repeatability]
\label{prop:affine-F}
Suppose that the probability assignment $P_\pm(\theta;a,b)$ satisfies Axioms~\ref{ax:normalization}--\ref{ax:regularity}, together with the repeatability Axiom~\ref{ax:repeatability}. Let $F$ denote the continuous function appearing in Lemma~\ref{lem:reduced-form}. Then $F$ is an affine function of its argument; that is, there exist real constants $\alpha,\beta$ such that
\begin{equation}
\label{eq:F-affine}
F(x) = \alpha x + \beta,\qquad \forall\,x\in[0,1].
\end{equation}
In conjunction with the symmetry condition $F(1-x)=1-F(x)$ from
\eqref{eq:F-symmetry}, this implies
\begin{equation}
\label{eq:F-identity}
F(x) = x,\qquad \forall\,x\in[0,1],
\end{equation}
so that Born’s rule is recovered:
\[
P_+(\theta;a,b)
=
\left|\,a\cos\frac{\theta}{2} + b\sin\frac{\theta}{2}\,\right|^2,
\qquad
P_-(\theta;a,b)
=
\left|\,a\sin\frac{\theta}{2} - b\cos\frac{\theta}{2}\,\right|^2.
\]
\end{proposition}

\begin{proof}[Proof sketch]
We provide only a schematic argument; a fully rigorous treatment can be obtained by suitably adapting standard operational derivations of Born’s rule to the present framework.

\medskip\noindent
\emph{Step 1: Mixtures and convexity.}
Consider an ensemble which, at the preparation stage, is a classical mixture of two pure states $|\psi_1\rangle$ and $|\psi_2\rangle$, with respective weights $p$ and $1-p$. By construction, the overall outcome statistics for the $\theta=0$ measurement form a convex combination of the statistics associated with the two subensembles:
\begin{equation}
\label{eq:mixture}
P_+^{\text{mix}}(0) = p\,P_+^{(1)}(0) + (1-p)\,P_+^{(2)}(0),
\end{equation}
where $P_+^{(k)}(0)$ denotes the probability of obtaining the outcome $+$ when the system is prepared in the state $|\psi_k\rangle$.

Moreover, any such classical mixture can be implemented operationally by randomly preparing either $|\psi_1\rangle$ or $|\psi_2\rangle$ with the prescribed probabilities, while leaving the measurement procedure fixed. The repeatability postulate (Axiom~\ref{ax:repeatability}) then entails that the asymptotic relative frequency of the outcome $+$ must coincide with the probability given by the mixture rule \eqref{eq:mixture}.

\medskip\noindent
\emph{Step 2: Application to the reduced form.}
We now specialise to the reference angle $\theta = 0$ and, invoking Lemma~\ref{lem:reduced-form}, write
\[
P_+(0;a,b) = F(|a|^2).
\]
Consider two pure states $|\psi_1\rangle$ and $|\psi_2\rangle$ whose overlaps with $|+\rangle_z$ have squared moduli $x_1$ and $x_2$, respectively. Then
\[
P_+^{(1)}(0) = F(x_1),\qquad P_+^{(2)}(0) = F(x_2).
\]
Form a classical mixture of these preparations with weights $p$ and $1-p$. By construction, the resulting ensemble yields an overall probability
\[
P_+^{\text{mix}}(0) = p\,F(x_1) + (1-p)\,F(x_2).
\]

Alternatively, one can prepare a single pure state whose overlap with $|+\rangle_z$ has squared modulus
\[
x := p x_1 + (1-p)x_2,
\]
and, by exploiting the repeatability of the measurement in the eigenbasis, arrange for the same long-run frequency of the outcome $+$. This operational equivalence imposes the constraint
\[
F(p x_1 + (1-p)x_2) = p\,F(x_1) + (1-p)\,F(x_2)
\]
for all $x_1,x_2\in[0,1]$ and all $p\in[0,1]$. Hence, $F$ must be an affine function on $[0,1]$.

\medskip\noindent
\emph{Step 3: Exploiting the symmetry condition.}
Let $F(x)=\alpha x+\beta$. Imposing the symmetry constraint $F(1-x)=1-F(x)$ for all $x\in[0,1]$ gives
\[
\alpha(1-x)+\beta = 1 - (\alpha x + \beta),
\]
which reduces to the linear system
\[
\alpha = 1,\qquad \beta = 0.
\]
Consequently, $F(x)=x$ on $[0,1]$, and the asserted expression for $P_\pm(\theta;a,b)$
follows directly from Lemma~\ref{lem:reduced-form}.
\end{proof}


\section{PDL interpretation and microscopic translation}
\label{sec:PDL-translation}

\subsection{PDL spin model and measurement interface}

Within the Projective Dynamic Logo (PDL) framework, physical systems are modelled as finite signed graphs determined by a coherence-optimisation principle.\cite{LaubscherLDP2026} The minimal $(4,6)$ closure assumes a dual function: it constitutes both the earliest logically admissible stationary regime singled out by the axiomatic structure and a canonical prototype for an effective two-level (spin-$\tfrac12$) system.\cite{LaubscherIntroPDL2026}

More precisely, we analyze the complete graph on four vertices with six signed edges, denoted \((4,6)\), equipped with a binary pulsation of edge signs that preserves triangular coherence at each update step. Within the class of admissible sign configurations, one can identify at least two distinct stationary motifs, \(\sigma^+\) and \(\sigma^-\), of identical coherence cost, which cannot be mapped onto one another by a trivial global sign inversion or by any relabeling of vertices that leaves all mixed-triangle couplings to external structures invariant.\cite{LaubscherBornPDL2026} These motifs are operationally distinguishable through their respective interaction profiles with a measurement apparatus and are therefore construed as logical representatives of the eigenstates \(\lvert + \rangle_z\) and \(\lvert - \rangle_z\).

A Stern–Gerlach-type measurement apparatus is modelled as a finite signed graph $G_{\mathrm{app}}$ endowed with an \emph{active surface} consisting of mixed triangles that couple the $(4,6)$ block to two distinguished interface branches.\cite{LaubscherBornPDL2026} More precisely, denoting by $V_e$ the set of vertices belonging to the $(4,6)$ block and by $I_+$, $I_-$ the interface vertices of the apparatus, one introduces two disjoint families of mixed triangles
\begin{align}
T_+ &= \bigl\{ (e_i,e_j,I_+)\,:\, e_i,e_j\in V_e,\ (e_i,e_j)\in E_e,\ 
(e_i,I_+),(e_j,I_+)\in E \bigr\},\\
T_- &= \bigl\{ (e_i,e_j,I_-)\,:\, e_i,e_j\in V_e,\ (e_i,e_j)\in E_e,\ 
(e_i,I_-),(e_j,I_-)\in E \bigr\},
\end{align}
where $E_e$ and $E$ denote, respectively, the edge set of the $(4,6)$ block and the edge set of the composite graph $G_{4,6}\cup G_{\mathrm{app}}$. Each triangle $\tau\in T_+\cup T_-$ is assigned a sign $\sigma_\tau(\sigma)$, determined by the edge-sign configuration $\sigma$, and a weight $\alpha_\tau(\theta,a,b)$ that encodes the dependence on the prepared state $|\psi\rangle=a|+\rangle_z + b|-\rangle_z$ as well as on the measurement direction $\theta$.

We then define a mixed-triangle coherence cost $\varepsilon_{\mathrm{mix}}(\sigma;\theta,a,b)$ by summing the weights of all violated triangles. This cost functional is used to specify an exponential weighting over global configurations via
\begin{equation}
w(\sigma;\theta,a,b)
=
\exp\bigl[-\lambda\,\varepsilon_{\mathrm{mix}}(\sigma;\theta,a,b)\bigr],
\qquad \lambda\gg 1.
\end{equation}
By partitioning the space of admissible configurations into two basins, $A$ and $B$, associated respectively with the predominance of mixed-triangle coherence in the $+$ and $-$ branches, and by normalising the total weights restricted to these basins, one obtains effective probabilities $P_\pm(\theta; a,b)$ which, in the limit of large $\lambda$, reproduce Born’s rule.\cite{LaubscherBornPDL2026} In the following, we examine how this PDL construction is related to the operational axioms introduced in Section~\ref{sec:operational-setting}.

\subsection{From PDL structure to the operational axioms}

We begin by demonstrating that, under mild symmetry assumptions imposed on the PDL measurement interface, the effective probabilities generated by the mixed-triangle scheme automatically satisfy Axioms \ref{ax:normalization}--\ref{ax:branch-symmetry}.

\begin{lemma}[PDL mixed-triangle scheme and operational axioms]
\label{lem:PDL-axioms}
Consider the PDL spin-measurement construction introduced above, involving two stationary motifs $\sigma^\pm$ on the $(4,6)$ block and two families of mixed triangles $T_\pm$ that determine the active surface of the apparatus. Suppose that the following conditions hold:
\begin{enumerate}[label=(\roman*)]
\item The mixed-triangle weights $\alpha_\tau(\theta,a,b)$ are non-negative for all $\tau\in T_+\cup T_-$, for all measurement directions $\theta$, and for all normalized parameter pairs $(a,b)$.
\item The apparatus is symmetric under exchange of the two outcome branches in the sense that there exists an involutive relabeling of $I_+$ and $I_-$ that maps $T_+$ onto $T_-$ and conversely, while preserving the internal coherence structure of $G_{\mathrm{app}}$.
\item A rotation of the measurement direction by an angle $\theta_0$ is represented, at the PDL level, by a relabeling of the mixed triangles compatible with the SU(2) action on the $(4,6)$ block, so that $\varepsilon_{\mathrm{mix}}(\sigma;\theta+\theta_0,a,b)$ is obtained from $\varepsilon_{\mathrm{mix}}(\sigma';\theta,a',b')$ for a suitably rotated configuration $(\sigma',a',b')$.
\item The basin weights
\[
W_A(\theta;a,b) = \sum_{\sigma\in A} w(\sigma;\theta,a,b),
\qquad
W_B(\theta;a,b) = \sum_{\sigma\in B} w(\sigma;\theta,a,b)
\]
are finite and strictly positive for all $\theta$ and all normalized $(a,b)$.
\end{enumerate}
Then the corresponding effective probabilities
\[
P_+(\theta;a,b)
=
\frac{W_A(\theta;a,b)}{W_A(\theta;a,b)+W_B(\theta;a,b)},\qquad
P_-(\theta;a,b)
=
\frac{W_B(\theta;a,b)}{W_A(\theta;a,b)+W_B(\theta;a,b)},
\]
satisfy Axioms
\ref{ax:normalization}--\ref{ax:branch-symmetry}.
\end{lemma}

\begin{proof}[Proof sketch]
Normalization and positivity (Axiom~\ref{ax:normalization}) follow directly from the definitions. By assumption, the weights $w(\sigma;\theta,a,b)$ are non-negative and not all simultaneously zero. Consequently, $W_A,W_B\ge 0$ and $W_A+W_B>0$, so that the ratios defining $P_\pm$ lie in the interval $[0,1]$ and sum to~1.

Global phase invariance (Axiom~\ref{ax:global-phase}) is inherited from the dependence of the mixed-triangle weights $\alpha_\tau(\theta,a,b)$ on the state amplitudes $(a,b)$. Since only phase-invariant combinations of $(a,b)$ enter the definition of $\varepsilon_{\mathrm{mix}}$, multiplication of $(a,b)$ by an overall phase leaves $w(\sigma;\theta,a,b)$ invariant, and therefore also $P_\pm(\theta;a,b)$.

Rotational covariance (Axiom~\ref{ax:rotation}) is encoded in assumption (iii). A rotation of the measurement axis by an angle $\theta_0$ induces a relabeling of triangles that is compatible with the SU(2) action on the $(4,6)$ block. As a consequence, $\varepsilon_{\mathrm{mix}}$ and thus $w$ transform in such a way that
\[
P_{\pm}(\theta+\theta_0;a,b) = P_{\pm}(\theta;a',b')
\]
for suitably rotated amplitudes $(a',b')$.

The eigenstate behavior at the reference angle (Axiom~\ref{ax:eigenstates}) is ensured by the construction of the stationary motifs $\sigma^\pm$. For $\theta=0$, a $(4,6)$ block in configuration $\sigma^+$ couples coherently only to mixed triangles in $T_+$, whereas $\sigma^-$ couples only to $T_-$. In the limit of large~$\lambda$, configurations in the ``incorrect'' basin are exponentially suppressed, implying $P_+(0;1,0)\to 1$ and $P_-(0;0,1)\to 1$.

Finally, branch symmetry (Axiom~\ref{ax:branch-symmetry}) follows from assumption (ii). Exchanging the labels $I_+$ and $I_-$ at the PDL level induces an interchange of $T_+$ and $T_-$, and therefore of the basins $A$ and $B$. Composing this exchange with a rotation by $\pi$ about an axis orthogonal to $z$ maps the probabilities $P_+(\theta;a,b)$ to $P_-(\theta+\pi;a,b)$, and conversely, thereby establishing the symmetry relations in Axiom~\ref{ax:branch-symmetry}.
\end{proof}

\section{Discussion and outlook}
\label{sec:discussion}

In this work, we have established, within an explicitly operational setting, a Gleason-type uniqueness theorem for spin-$\tfrac12$ measurements that is tailored to the PDL framework. Starting from axioms of normalization, global phase invariance, rotational covariance, branch symmetry, and regularity, we first demonstrated that any admissible probability assignment $P_\pm(\theta;a,b)$ must depend on the prepared state $|\psi\rangle = a|+\rangle_z + b|-\rangle_z$ and on the measurement direction $\theta$ only through the scalar quantity $\bigl|a\cos(\theta/2)+b\sin(\theta/2)\bigr|^2$, up to a continuous function $F:[0,1]\to[0,1]$ satisfying $F(1-x)=1-F(x)$. By incorporating a natural repeatability postulate and exploiting the behaviour of classical mixtures, we derived Jensen’s functional equation for $F$ and concluded that $F(x)$ must be affine and, by symmetry, equal to the identity. Born’s rule thus arises as the \emph{unique} probability assignment compatible with our operational axioms in the spin-$\tfrac12$ setting.

On the PDL side, we have re-examined the mixed-triangle measurement scheme introduced in Ref.~\cite{LaubscherBornPDL2026}, in which a $(4,6)$ block encoding a logical spin doublet is coupled to a Stern–Gerlach-type apparatus through a finite active surface. Under mild symmetry assumptions on the apparatus—specifically, an exchange symmetry between the two outcome branches and compatibility with the SU(2) action on the $(4,6)$ structure—we rigorously established that the resulting effective probabilities $P_\pm(\theta;a,b)$ satisfy our operational axioms. When combined with the uniqueness theorem, this result implies that the PDL scheme does not merely \emph{realize} Born’s rule for spin-$\tfrac12$ measurements, but in fact lies within a class of models in which Born’s rule is \emph{uniquely} enforced by the imposed structural constraints.

These findings constitute a first explicit bridge between Gleason-type derivations and the discrete relational substrate underlying PDL, and they motivate several avenues for further research. At the operational level, a natural extension is to move beyond two-outcome observables and to determine to what extent higher-dimensional spin systems or general POVMs can be accommodated within an analogous axiomatic setting. At the level of PDL itself, one would like to relax some of the symmetry assumptions imposed on the apparatus, to construct fully dynamical models of the measurement process, and to investigate whether comparable uniqueness results can be obtained for other classes of PDL observables. Ultimately, a comprehensive Gleason-type theory for PDL would aim to characterise, in purely combinatorial terms, the full set of probability assignments compatible with the axioms of the framework and to isolate Born’s rule as a structurally forced outcome rather than as a contingent modelling choice.

From a more conceptual perspective, the present result supports the claim that the quantum probability structure associated with spin-$\tfrac12$ systems is not a freely tunable component within PDL, but rather a structurally compelled consequence of the underlying relational axioms together with the assumed operational symmetries. Once the discrete PDL substrate and the measurement principles of Section~\ref{sec:operational-setting} are fixed, Born’s rule no longer appears as an independent postulate; instead, it emerges as an unavoidable implication of the combined structural and operational constraints. While we refrain from drawing broader metaphysical conclusions concerning determinism or free will, this viewpoint suggests that key aspects of quantum phenomenology are subject to stringent structural restrictions—and may even be rendered effectively inevitable—within a suitably configured relational framework.

\bibliographystyle{unsrt}
\bibliography{PDL_Gleason}

\end{document}
